\section{Reaction--diffusion: vì sao mức quan tâm lan truyền?}
\label{sec:reaction-diffusion}

\subsection{Hai bể nước và một ý tưởng}

Trong hóa học, một chất khuếch tán từ nơi có nồng độ cao sang nơi có nồng độ thấp. D-Tracker
mượn hình ảnh đó cho mức quan tâm: nếu một từ khóa đang được chú ý mạnh ở nhóm địa điểm $j'$,
sự chú ý có thể chảy sang nhóm $j$ đang thấp hơn. Phương trình reaction--diffusion tổng quát mà
paper dẫn ra có dạng
\begin{equation}
  \frac{\partial u_i}{\partial t}
  =f(u_i)+\sum_j D\,(u_j-u_i),
  \label{eq:rd-background}
\end{equation}
trong đó $u_i$ là nồng độ tại vị trí $i$ và $D$ là độ mạnh khuếch tán. Số hạng $f(u_i)$ được gọi
là \textbf{reaction}: nó mô tả thay đổi cục bộ tại chỗ. Số hạng $\sum_j D(u_j-u_i)$ là
\textbf{diffusion}: nó chỉ khác 0 khi các vị trí có mức khác nhau, và dấu của nó cho biết dòng
chảy vào hay ra.

\begin{center}
\begin{tikzpicture}[
  >=Latex,
  place/.style={circle, draw=DTNavy, thick, minimum size=18mm, fill=DTMist,
                align=center, font=\sffamily\small},
  flow/.style={->, very thick, DTTeal},
  label/.style={font=\small, align=center}
]
  \node[place] (b) at (0,0) {nhóm $j'$\\$w_{ij'}=4$};
  \node[place] (a) at (4.2,0) {nhóm $j$\\$w_{ij}=1$};
  \draw[flow] (b) to[bend left=12] node[above,label] {dòng vào vì $4>1$} (a);
  \draw[flow, DTClaret] (a) to[bend left=12] node[below,label] {dòng ra nếu $1>4$} (b);
\end{tikzpicture}
\end{center}

\subsection{Phiên bản latent của D-Tracker}

D-Tracker không áp đặt phương trình lên từng cặp từ khóa--địa điểm. Nó áp đặt lên \emph{lõi ẩn}
gồm $d_k\cdot d_l$ dynamics. Với $i=1,\ldots,d_k$ là nhóm từ khóa và $j=1,\ldots,d_l$ là nhóm
địa điểm, paper viết
\begin{equation}
  \frac{dw_{ij}}{dt}
  =a_{ij}w_{ij}
   +\sum_{j'=1}^{d_l} d_{ijj'}\,(w_{ij'}-w_{ij}),
  \qquad 1\leq i\leq d_k,\;1\leq j\leq d_l.
  \label{eq:rd-latent}
\end{equation}
Hai tham số mới cần đọc đúng:
\begin{itemize}
  \item $a_{ij}$ là \textbf{tốc độ tăng trưởng}. Nếu $a_{ij}>0$, dynamics $w_{ij}$ có xu hướng
        tăng; nếu $a_{ij}<0$, nó có xu hướng giảm. Đây là ``reaction''.
  \item $d_{ijj'}$ là \textbf{độ mạnh khuếch tán} từ nhóm địa điểm $j'$ sang nhóm $j$. Trong paper,
        $\mathcal{D}\in\mathbb{R}_{+}^{d_k\times d_l\times d_l}$ được mô tả là không âm. Dấu của
        dòng chảy thực tế còn phụ thuộc vào $(w_{ij'}-w_{ij})$: nếu $w_{ij'}>w_{ij}$ thì $j$ nhận
        thêm; nếu nhỏ hơn thì $j$ mất bớt.
\end{itemize}

Ở đây $w_{ij}$ là mức hoạt động của dynamics thuộc nhóm từ khóa $i$ và nhóm địa điểm $j$. Các
dynamics khác nhóm từ khóa không khuếch tán trực tiếp trong số hạng này; chúng tương tác qua lõi
và qua phép chiếu đa tuyến tính ở Chương~\ref{sec:model}. Toàn bộ hệ có $d_k\cdot d_l$ phương
trình, tham số gồm ma trận $\mathbf{A}\in\mathbb{R}^{d_k\times d_l}$ và tensor
$\mathcal{D}\in\mathbb{R}_{+}^{d_k\times d_l\times d_l}$.

\begin{keypoint}
$w>0$ giúp đọc ``mức hoạt động''; $a$ nói tăng hay giảm; $(w_{j'}-w_j)$ nói hướng chảy; $d$ nói
dòng chảy mạnh đến đâu. Đây là lý do mô hình có thể diễn giải được, chứ không chỉ là một hộp đen
dự báo.
\end{keypoint}

\subsection{Một bước số cụ thể}

Giả sử một dynamics ở nhóm $j$ có $w_j=1$, một dynamics cùng nhóm từ khóa ở nhóm $j'$ có
$w_{j'}=4$, $a_j=0.1$ và $d_{jj'}=0.5$. Bỏ qua các nhóm khác, đạo hàm tại thời điểm này là
\[
  \frac{dw_j}{dt}=0.1\cdot1+0.5\cdot(4-1)=0.1+1.5=1.6.
\]
Nếu xấp xỉ Euler với bước $\Delta t=1$, giá trị kế tiếp xấp xỉ $1+1.6=2.6$. Dynamics đang tăng
vì vừa có growth vừa nhận dòng từ nơi cao hơn. Nếu đổi $w_{j'}=0$, số hạng khuếch tán thành
$0.5(0-1)=-0.5$: cùng một độ mạnh $d$ nhưng dòng chảy đổi chiều.

\begin{plainwords}
Đừng học thuộc phương trình. Hãy đọc ba câu: ``tự nó tăng/giảm'', ``nó bị kéo về phía nhóm khác'',
và ``mức chênh càng lớn thì lực kéo càng mạnh''. Tất cả nằm trong hai số hạng của
Công thức~\eqref{eq:rd-latent}.
\end{plainwords}

\subsection{Liên hệ với dữ liệu quan sát}

Lõi $w_{ij}$ là ẩn. Để nối nó với chuỗi quan sát $X_{uv}$ của từ khóa thật $u$ và địa điểm thật
$v$, paper lấy đạo hàm của phép chiếu và thu được
\begin{equation}
  \frac{dX_{uv}}{dt}
  =\sum_{i=1}^{d_k}\sum_{j=1}^{d_l}
   W^{(\mathrm{key})}_{iu}W^{(\mathrm{loc})}_{jv}
   \underbrace{\left(
     a_{ij}w_{ij}+\sum_{j'=1}^{d_l}d_{ijj'}(w_{ij'}-w_{ij})
   \right)}_{\text{hệ reaction--diffusion ẩn}}.
  \label{eq:observed-derivative}
\end{equation}
Đọc từ phải sang trái: hệ ẩn tạo ra $d_k d_l$ động học; các trọng số không âm
$W^{(\mathrm{key})}_{iu}W^{(\mathrm{loc})}_{jv}$ quyết định dynamics nào đóng góp vào cặp
$(u,v)$. Vì trọng số không âm, phép chiếu không đảo dấu: một dynamics tăng vẫn đóng góp dương
vào chuỗi quan sát tương ứng. Đây là cầu nối giữa ``lõi nhỏ'' và ``chuỗi thật''.

\subsection{Diễn giải không đồng nghĩa với nhân quả}

Giả định reaction--diffusion nói rằng mức chênh có thể tạo ra dòng chảy. Mô hình học các tham số
$a,d$ và hệ số nhóm để khớp dữ liệu; nó \emph{không} chứng minh rằng sự tăng ở Mỹ gây ra sự tăng
ở Việt Nam. Một yếu tố thứ ba (ra mắt sản phẩm toàn cầu, tin tức, mạng xã hội) có thể tạo ra cùng
một pattern trễ. Vì vậy, hãy đọc kết quả là \textbf{giả thuyết lan truyền có cấu trúc}, cần kiểm
tra bằng kiến thức miền và thí nghiệm bổ sung.

\begin{commonmistake}
``$d_{ijj'}>0$ nghĩa là $j'$ luôn truyền sang $j$'' là chưa đủ. Dòng chảy còn phụ thuộc vào
$w_{ij'}-w_{ij}$; nếu $w_{ij'}<w_{ij}$, số hạng này âm. Paper gọi $d$ là độ mạnh, không phải một
chiều cố định. Tương tự, ``$a>0$'' chỉ nói xu hướng cục bộ, không đảm bảo tăng mãi mãi.
\end{commonmistake}

\begin{selfcheck}
\begin{enumerate}
  \item Với $w_j=2,w_{j'}=5,a_j=0,d_{jj'}=0.2$, đạo hàm bằng bao nhiêu?
  \item Nếu $a_{ij}<0$ và dòng khuếch tán dương, dynamics có chắc chắn giảm không? Vì sao?
  \item Nêu một tình huống dữ liệu có độ trễ giữa hai nơi nhưng không phải lan truyền nhân quả.
\end{enumerate}
\end{selfcheck}
